% !TeX root = ../main.tex

\section{Методы}

\subsection*{Модели генерации BOLD-сигнала}

Рассматриваются две модели формирования BOLD-сигнала: P-DCM~\cite{havlicek2015} и классическая модель Balloon--Windkessel с сопряжением типа Фристона (далее BW). Обе модели представляют собой каскад: общий линейный блок возбуждения--торможения (E/I-блок) преобразует вход в траекторию нейрональной активности, сопряжение переводит ее в изменение кровотока, сосудистый блок описывает динамику венозного объема и дезоксигемоглобина, а уравнение наблюдения задает BOLD-сигнал. E/I-блок у обеих моделей один и тот же, поэтому при одинаковом входе траектории нейрональной активности совпадают, а вся изучаемая нелинейность сосредоточена в сосудистом блоке и уравнении наблюдения.

Состоянием покоя называется устойчивое равновесие системы при нулевом входе. Кровоток $f$, объем $v$ и количество дезоксигемоглобина $q$ нормированы на значения покоя и в покое равны единице; отклонения $x_E$, $x_I$ и $z$ в покое равны нулю. Переменная $z$ описывает нейроваскулярное сопряжение в модели и не является отдельно регистрируемым физиологическим сигналом. Общий нейрональный блок с входом $u(t)$, заданным в~\eqref{eq:input}:
\begin{equation}\label{eq:a1}
\dot x_E=-\sigma x_E-\mu x_I+u,\qquad
\dot x_I=\lambda(x_E-x_I),
\qquad
\sigma=0{,}5,\quad \mu=0{,}4,\quad \lambda=0{,}2~\text{с}^{-1}.
\end{equation}
Сопряжение у моделей различается. Для BW оно задается уравнениями $\dot z=x_E-\kappa z-\gamma(f-1)$, $\dot f=z$, где $\kappa=0{,}64$, $\gamma=0{,}32$: кровоток через обратную связь действует на переменную нейроваскулярного сопряжения $z$. Для P-DCM $\dot z=x_E-\varphi z$, $\dot f=\phi z-\chi(f-1)$, где $\varphi=0{,}6$ --- затухание $z$, $\phi=1{,}5$ --- усиление кровотока, $\chi=0{,}6$ --- затухание кровотока; обозначения следуют~\cite{havlicek2015}.

Сосудистый блок описывает баланс притока и оттока крови, образование дезоксигемоглобина при извлечении кислорода и его удаление с оттоком. Для обеих моделей
\begin{equation}\label{eq:a2}
\tau_0\dot v=f-f_{\mathrm{out}},\qquad
\tau_0\dot q=\frac{fE(f)}{E_0}-f_{\mathrm{out}}\frac qv,
\qquad
E(f)=1-(1-E_0)^{1/f}.
\end{equation}
В BW $f_{\mathrm{out}}=v^{1/\alpha}$; в вязкоупругом варианте P-DCM $f_{\mathrm{out}}=v^{1/\alpha}+\tau_v\dot v$, откуда $\dot v=(f-v^{1/\alpha})/(\tau_0+\tau_v)$. В основной паре моделей использованы время среднего транзита крови $\tau_0=2$~с, показатель Грабба $\alpha=0{,}32$ и исходная доля извлечения кислорода $E_0=0{,}4$; постоянная вязкоупругого оттока $\tau_v=4$~с в P-DCM и $\tau_v=0$ в BW. Нелинейность переноса $q/v$ сохранена. Общее наблюдение BOLD в долях исходного сигнала:
\begin{equation}\label{eq:a3}
\begin{split}
y&=V_0\left[c_1(1-q)+c_2(1-q/v)+c_3(1-v)\right],\\
V_0&=0{,}04,\quad c_1=4{,}15896,\quad c_2=0{,}6204,\quad c_3=0{,}53.
\end{split}
\end{equation}
Здесь $q$ --- относительное количество дезоксигемоглобина, $q/v$ --- его относительная концентрация, $V_0$ --- исходная венозная объемная доля; коэффициенты наблюдения фиксированы для поля 3~Тл и не подбирались по генератору или тестовому ряду.

Источниками нелинейности служат степенной закон оттока $v^{1/\alpha}$, зависимость извлечения кислорода из крови от кровотока, перенос дезоксигемоглобина через отношение $q/v$ и уравнение наблюдения; линейные E/I-блок и сопряжение собственных нелинейных членов не добавляют. При переходе от P-DCM к BW меняются сразу сопряжение с кровотоком и закон оттока, поэтому различия результатов двух генераторов нельзя приписать одному механизму. Обе модели локальны и не описывают пространственную структуру сосудистой сети.

\subsection*{Входной сигнал и заданные амплитуды}

Временная структура двух условий задается бинарными функциями $s_A(t)$ и $s_B(t)$, принимающими значение 1 во время соответствующего предъявления и 0 вне его:
\begin{equation}\label{eq:input}
u(t)=A_0\,[a_A s_A(t)+a_B s_B(t)], \qquad
\Delta=a_A-a_B,
\end{equation}
где $a_A$ и $a_B$ --- заданные безразмерные амплитуды, а $\Delta$ --- заданная разность входных амплитуд в модельной шкале; с ней сравнивается оценка контраста ОЛМ. Общий множитель $A_0\approx0{,}88$ зафиксирован по номинальной P-DCM, то есть модели с исходными параметрами: единичная амплитуда секундного события дает единичный пиковый прирост относительного кровотока. Для BW используется тот же вход. Значение $a=0{,}2$ задано в выбранной модельной шкале и не является абсолютной мерой нейрональной активности человека. Общий линейный E/I-блок исключает нелинейную нейрональную адаптацию, поэтому вся регистрируемая нелинейность имеет гемодинамическое происхождение.

\subsection*{Оценка амплитуд и контраста}

Стандартная оценка выполнялась общей линейной моделью (ОЛМ) с канонической функцией гемодинамического ответа пакета SPM (далее $C$) --- разностью двух гамма-функций с формами компонент 6 и 16, единичными масштабами, относительным весом отрицательной компоненты $1/6$ и окном 32~с. Масштаб $g_0\approx0{,}375$ согласует пик линейного ответа $C$ на секундное событие с номинальной P-DCM, поэтому ошибка формы не включает тривиальное различие высоты. Форма, задержка и усиление $C$ по тестовым данным не оценивались.

Для каждого экспериментального условия временной ряд предъявления стимулов сворачивали с канонической HRF ($C$). Полученные регрессоры и константу включали в ОЛМ. Оценки выполнялись обобщенным методом наименьших квадратов с известной ковариацией шума. Основными величинами были оценки амплитуд $\hat\beta_{C,A}$, $\hat\beta_{C,B}$ и контраста $\hat\Delta_C=\hat\beta_{C,A}-\hat\beta_{C,B}$, сравниваемые с заданными $a_A$, $a_B$ и $\Delta$.

Наряду со стандартной оценкой для каждого генератора из его уравнений вычислялось линейное приближение около покоя (далее $L_1$). Регрессия с $L_1$ служит только вычислительным контролем: она устраняет ошибку выбора линейной формы и оставляет ошибку конечных амплитуд. На реальных данных $L_1$ заранее неизвестно.

Связь расхождения во временной области с ошибками оценок задается разложением сигнала. После добавления шума
\begin{equation}\label{eq:decomp}
\mathbf y=\mathbf X_1\mathbf a+\mathbf r_{\mathrm{nl}}+\boldsymbol\varepsilon,
\end{equation}
где $\mathbf a=(a_A,a_B)^\top$ --- заданные амплитуды, $\mathbf X_1$ --- регрессоры $L_1$, $\mathbf r_{\mathrm{nl}}$ --- нелинейный остаток, то есть разность полного ответа и линейной составляющей, $\boldsymbol\varepsilon$ --- шум. Все величины берутся после одинаковой обработки: выбора отсчетов, учета ковариации и центрирования. Обозначим через $\mathbf X_C$ матрицу регрессоров с $C$, через $\mathbf L_C$ --- оператор оценки амплитуд, а через $\ell_C=c^\top\mathbf L_C$ --- строку для контраста $c=(1,-1)^\top$. Тогда
\begin{equation}\label{eq:errdec}
\hat\Delta_C-\Delta=
\ell_C(\mathbf X_1-\mathbf X_C)\mathbf a+
\ell_C\mathbf r_{\mathrm{nl}}+
\ell_C\boldsymbol\varepsilon.
\end{equation}
Слагаемые описывают вклады несовпадения формы и масштаба HRF, нелинейного остатка и шума.

Решающую роль играет проекция $\ell_C\mathbf r_{\mathrm{nl}}$. Нелинейный остаток, похожий на регрессор одного условия, целиком переходит в соответствующий коэффициент, а одинаковый для обоих условий вклад при вычитании в основном сокращается. Поэтому заметное расхождение временного ряда может слабо изменить контраст, а небольшое расхождение может изменить его сильно, и качество описания ряда само по себе не является мерой точности оценок. При $L_1$ ошибка формы исчезает, и в бесшумном случае ошибка контраста определяется только проекцией нелинейного остатка. Сам остаток $\mathbf r_{\mathrm{nl}}$ не совпадает с остатком подгонки ОЛМ, поскольку часть нелинейности уже перешла в коэффициенты.

\subsection*{Разложение Вольтерра до второго порядка}

Нелинейный остаток анализировался разложением генератора в окрестности покоя:
\begin{equation}\label{eq:volterra}
\mathcal B[\epsilon u]
=\epsilon H_1[u]+\epsilon^2Q[u,u]+O(\epsilon^3),
\end{equation}
где $\mathcal B$ --- полный генератор BOLD, $H_1$ --- линейный оператор, совпадающий с $L_1$, $Q$ --- квадратичный оператор, $\epsilon$ --- масштаб входа. Квадратичный член уменьшается вчетверо при двукратном уменьшении входа. Ряды Вольтерра применяются к нелинейным системам с памятью~\cite{volterra1930,schetzen1980}, включая физиологические~\cite{marmarelis2004,hunter1986}; для систем с затухающей памятью существуют результаты об аппроксимации такими рядами~\cite{boyd1985}. Обычно ядра оценивают по наблюдаемым данным~\cite{pintelon2012,schoukens2015}; для BOLD-сигнала ядра Вольтерра оценивались по данным фМРТ~\cite{friston1998} и рассчитывались для модели Balloon~\cite{friston2000}. Здесь квадратичное ядро получено аналитически из уравнений обоих генераторов.

Квадратичная часть содержит вклады одиночных событий и пар событий. Первые описывают отклонение конечного ответа от линейного масштабирования, вторые --- нарушение суперпозиции из-за влияния предыдущих событий, возникающее и внутри одного условия, и между условиями. Само разложение не идентифицирует конкретный сосудистый механизм: сходные зависимости вход--выход возможны при разных сочетаниях динамических и статических преобразований.

Обозначим шестимерное состояние $x=(z,f,v,q,x_E,x_I)^\top$, покой $x_0=(0,1,1,1,0,0)^\top$, систему $\dot x=F(x)+\mathbf b u$ и наблюдение $y=G(x)$, где $\mathbf b=(0,0,0,0,1,0)^\top$. Пусть $J=DF(x_0)$ --- гурвицева матрица Якоби (все собственные значения лежат в левой полуплоскости), $d=DG(x_0)$ --- строка производных наблюдения, $H_F=D^2F(x_0)$ и $H_G=D^2G(x_0)$ --- гессианы. При гладкости в окрестности покоя и старте из состояния $x_0$ точное первое ядро имеет вид
\begin{equation}\label{eq:h1}
h_1(\tau)=d\,e^{J\tau}\mathbf b,\quad \tau\ge0;\qquad h_1(\tau)=0,\quad\tau<0.
\end{equation}
Непрерывное симметричное ядро второго порядка определяется суммой
\begin{equation}\label{eq:k2}
k_2(\tau_1,\tau_2)=k_{2,\mathrm{field}}(\tau_1,\tau_2)+k_{2,\mathrm{obs}}(\tau_1,\tau_2),
\end{equation}
\begin{equation}\label{eq:k2field}
k_{2,\mathrm{field}}(\tau_1,\tau_2)=\frac12 d\int_0^{\min(\tau_1,\tau_2)}e^{Js}H_F[e^{J(\tau_1-s)}\mathbf b,\,e^{J(\tau_2-s)}\mathbf b]\,ds,
\end{equation}
\begin{equation}\label{eq:k2obs}
k_{2,\mathrm{obs}}(\tau_1,\tau_2)=\frac12[e^{J\tau_1}\mathbf b]^\top H_G[e^{J\tau_2}\mathbf b].
\end{equation}
При отрицательном хотя бы одном лаге $k_2=0$. Множитель $1/2$ уже включен в ядро, поэтому оператор
\begin{equation}\label{eq:Q}
Q[u_1,u_2](t)=\int_0^\infty\int_0^\infty k_2(\tau_1,\tau_2)\,u_1(t-\tau_1)\,u_2(t-\tau_2)\,d\tau_1d\tau_2
\end{equation}
не требует дополнительного внешнего множителя. Гессианы разделяют четыре источника квадратичной нелинейности: степенной закон оттока $v^{1/\alpha}$ в $\dot v$, нелинейное извлечение кислорода $g(f)=f[1-(1-E_0)^{1/f}]/E_0$ в $\dot q$, конвективный перенос дезоксигемоглобина $f_{\mathrm{out}}q/v$ и нелинейное уравнение наблюдения BOLD. Линейные E/I-блок и сопряжение с притоком собственных квадратичных членов не добавляют. Все производные вычислялись автоматическим дифференцированием уравнений генератора.

Для разложения по событиям положим $u=\sum_{i\in A}u_i+\sum_{j\in B}u_j$, включив амплитуду и масштаб в каждое $u_i$. Квадратичная часть распадается на пять компонентов:
\begin{equation}\label{eq:five}
y_2=S_A+S_B+P_{AA}+P_{BB}+P_{AB},
\end{equation}
\begin{equation}\label{eq:fiveb}
\begin{split}
S_A&=\sum_{i\in A}Q[u_i,u_i],\qquad S_B=\sum_{i\in B}Q[u_i,u_i],\\
P_{AA}&=2\sum_{i<j,\;i,j\in A}Q[u_i,u_j],
\end{split}
\end{equation}
и аналогично $P_{BB}$, $P_{AB}=2\sum_{i\in A,j\in B}Q[u_i,u_j]$. Полный нелинейный остаток во временной области есть $r_{\mathrm{nl}}=y_2+r_{\ge3}$, где $r_{\ge3}$ объединяет вклады третьего и высших порядков. Проекция каждого слагаемого на тот же дизайн дает его вклад в ошибку контраста.

Из полного квадратичного ядра выделена статическая проекция --- винеровский член $\kappa_0\,h_1\otimes h_1$:
\begin{equation}\label{eq:kappa}
\kappa_0=\frac{\langle k_2,h_1\otimes h_1\rangle}{\|h_1\otimes h_1\|^2},\qquad
k_{2,\perp}=k_2-\kappa_0\, h_1\otimes h_1,
\end{equation}
где скалярное произведение берется в норме $L^2$ на квадрате лагов $[0,T]^2$, $T=64$~с.

\subsection*{Аналитические поправки к оценкам канонической модели}

Обе поправки используют только уже полученные оценки $\hat{\boldsymbol\beta}_C$ и известную последовательность предъявления стимулов; новое обращение к BOLD-ряду, оптимизация и итерации отсутствуют. По известной последовательности и $\hat{\boldsymbol\beta}_C$ строится оцененный вход $\hat u_C(t)=A_0[\hat\beta_{C,A}s_A(t)+\hat\beta_{C,B}s_B(t)]$. По заранее заданной номинальной модели вычисляется нелинейный вклад $b_0[\hat u_C]$, проекция которого однократно вычитается из оценок коэффициентов:
\begin{equation}\label{eq:corr}
\hat{\boldsymbol\beta}_{\mathrm{corr}}
=\hat{\boldsymbol\beta}_C-
\mathbf L_C\, b_0[\hat u_C],
\qquad
\hat\Delta_{\mathrm{corr}}
=c^\top\hat{\boldsymbol\beta}_{\mathrm{corr}}.
\end{equation}
Идейно схема связана с обращением ряда возмущений: однократное вычитание оцененного члена второго порядка соответствует первому шагу обратного ряда рассеяния~\cite{moskow2008}. Оценка с $C$ возвращается без поправки, если $\max|\hat\beta_C|>0{,}45$, если отношение минимального сингулярного числа обработанного дизайна к максимальному меньше $10^{-6}$ или если результат не является конечным числом. Поскольку $\hat u_C$ содержит ошибки исходной регрессии и шум, реализуемая поправка отличается от теоретического расчета с входными амплитудами.

Поправки различаются только оператором квадратичного вклада. Полная поправка Вольтерра (Full Volterra, далее FV) использует $b_0[u]=Q_0[u,u]$, полный квадратичный оператор номинальной P-DCM. Статическая поправка винеровского типа (далее FW) заменяет его на $b_0[\hat u_C]=\kappa_0\,(H_0[\hat u_C])^2$: номинальный линейный оператор $H_0$ с ядром $h_1$ фильтрует оцененный вход, результат возводится в квадрат, а коэффициент $\kappa_0$ вычислен проекцией~\eqref{eq:kappa}. Параметры обеих поправок были зафиксированы до всех испытаний.

\subsection*{Показатели нелинейности и ошибки оценок}

Основными метриками были ошибки оценок $\beta$ и контраста относительно заданных амплитуд и $\Delta$. Для отдельной реализации знаковая ошибка контраста равна $e=\hat\Delta-\Delta$. При ненулевой разности используется относительная ошибка $100e/\Delta$, а для коэффициента при ненулевой амплитуде --- $100e_{\beta_j}/a_j$, $j\in\{A,B\}$. При вычислении относительной среднеквадратичной ошибки (root mean square error, RMSE) ошибки нормируются на заданную разность амплитуд или амплитуду до возведения в квадрат и агрегирования; при нулевой разности относительная ошибка не определена.

В бесшумных сериях источником изменчивости служит только последовательность предъявления стимулов, поэтому ошибки описываются медианой с диапазоном min--max. В шумных условиях реализации различаются последовательностью и шумом, поэтому дополнительно используется RMSE, учитывающая и смещение, и разброс; средняя знаковая ошибка по восьми повторам дает приближение условного смещения, то есть математического ожидания ошибки при фиксированных параметрах и последовательности стимулов.

Для характеризации нелинейности ответов на одиночные и парные предъявления стимулов использовались две величины. Взаимодействие пары событий определялось как $I(t)=y[u_1+u_2](t)-y[u_1](t)-y[u_2](t)$, где второе и третье слагаемые --- ответы на изолированные события с теми же сдвигами; в линейной системе $I(t)\equiv0$, а знак интеграла $S_I=\int I(t)\,dt$ показывает, ниже или выше суммы ответ пары. Отклонение одиночного ответа от пропорциональности оценивалось нормированным остатком $[y(t;a)-a\,y_1(t)]/a^2$, где $y_1$ --- линейный ответ на единичную амплитуду. Совпадение кривых, нормированных на $a^2$, служит признаком локального квадратичного масштабирования, расхождение указывает на вклад высших порядков. Влияние этих вкладов на оценки $\beta$ и контраст определяется их проекцией на дизайн.
